\documentclass[letterpaper,11pt]{article}
\usepackage{graphicx}
\usepackage{geometry}
\linespread{1.2}                        
\setlength{\oddsidemargin}{0.0in}            
\setlength{\textwidth}{6.5in}    
\usepackage[utf8]{inputenc} % Asegúrate de usar utf8
\usepackage[T1]{fontenc} % Asegúrate de usar T1 font encoding      
\setlength{\topmargin}{-.6in}            
\setlength{\textheight}{9in}
\usepackage{latexsym} 
\usepackage{amssymb} 
\usepackage{amsmath}
\usepackage{amsxtra}
\usepackage[mathcal]{eucal} 
\usepackage{enumerate}
\usepackage{hyperref}
\usepackage{rotating}
\usepackage{caption}
\usepackage{lscape}
\usepackage{paralist} 
\usepackage{indentfirst}
\usepackage[dvipsnames,svgnames]{xcolor} 
\usepackage{setspace}
\usepackage{multicol} 
\usepackage{float}

\definecolor{codegreen}{rgb}{0,0.6,0}
\definecolor{codegray}{rgb}{0.5,0.5,0.5}
\definecolor{codepurple}{cmyk}{0.65, 1, 0, 0.35}
\definecolor{backcolour}{rgb}{0.95,0.95,0.92}

\usepackage{listings}
\usepackage{xcolor}

\lstset{
    language=R, % lenguaje de programación
    basicstyle=\small\ttfamily, % estilo básico
    commentstyle=\color{gray}, % estilo para los comentarios
    keywordstyle=\color{blue}, % estilo para las palabras clave
    stringstyle=\color{black}, % estilo para las cadenas de texto
    showstringspaces=false, % no mostrar espacios en cadenas de texto
    tabsize=4, % tamaño de la tabulación
    breaklines=true, % permitir el salto de línea automático
    postbreak=\mbox{\textcolor{red}{$\hookrightarrow$}\space}, % símbolo al final de una línea partida
    numbers=left, % mostrar números de línea
    numberstyle=\tiny\color{gray}, % estilo para los números de línea
    frame=single, % mostrar un borde alrededor del código
    backgroundcolor=\color{gray!10}, % color de fondo
}

\begin{document}

\begin{center}
    {\Large  Taller 1: Econometría }

Profesora:    Jocelyn Tapia Stefanoni\\
Ayudantes:   Matías Balboa y   Barbara Buffo
\end{center}

\textbf{Instrucciones:} 

\begin{itemize}
    \item Usted dispone de 60 minutos  para intentar los 100 puntos del taller.
    \item Cuide la presentación y redacción de sus respuestas.
    \item Puede utilizar su computador y los apuntes tanto de clase como de ayudantía.
\item Descargue sus respuestas en un archivo con la extensión .ipynb. Debe enviar el enlace al notebook en Google Colab junto con el archivo .ipynb adjunto por correo electrónico a la profesora (jocelyn.tapias@usm.cl), con copia a los ayudantes (econometriausm482@gmail.com). 
    \end{itemize}

Para este taller, utilice los datos del archivo \texttt{wage2} del paquete wooldridge. 


\begin{enumerate}
\item (10 puntos) Describa la base de datos, identificando su origen y las variables que contiene. 
    \item (10 puntos) Calcule y presente estadísticos descriptivos relevantes para todas las variables contenidas en la base de datos.

\item (10 puntos) Obtenga e interprete la matriz de correlaciones entre las variables incluidas en la base de datos.

\item (10 puntos) Estime una regresión lineal simple del \texttt{IQ} sobre \texttt{educ}, obteniendo el coeficiente de la pendiente, que denotará como $\tilde{\delta}_1$.

\item (10 puntos) Estime una regresión lineal simple del logaritmo del salario $\log(\texttt{wage})$ sobre \texttt{educ}, obteniendo el coeficiente de la pendiente, que denotará como $\tilde{\beta}_1$.

\item (10 puntos) Estime una regresión lineal múltiple del logaritmo del salario $\log(\texttt{wage})$ sobre \texttt{educ} y \texttt{IQ}, obteniendo los coeficientes de las pendientes, que denotará como $\hat{\beta}_1$ y $\hat{\beta}_2$, respectivamente.

\item (14 puntos) Verifique empíricamente la siguiente relación:
\[
\tilde{\beta}_1 \approx \hat{\beta}_1 + \hat{\beta}_2 \cdot \tilde{\delta}_1
\]
donde $\tilde{\delta}_1$ es obtenido del punto (4), $\tilde{\beta}_1$ del punto (5), y $\hat{\beta}_1$ y $\hat{\beta}_2$ del punto (6).

\item (16 puntos) Determine e interprete la variabilidad explicada y la variabilidad no explicada 
 por el modelo múltiple estimado en el punto (6).

 \item (10 puntos) Determine el error estándar de los estimadores de MCO del modelo del punto (6). 
%\item (22 puntos) Sugiera y justifique claramente qué variables adicionales podría agregar al modelo para mejorar su especificación. Incorpórelas, estime nuevamente el modelo, y compare e interprete sus resultados en relación al modelo original estimado en el punto (6).
\end{enumerate}
\end{document}